\section{Control software}\label{sec:control-software}
% Draft by Claude, edit freely. Overall functionality and why, not every button.
% Keep adding to this section until the semester ends.

The Raspberry Pi runs a small Python program that controls the device. It runs the sampling
cycle on a timer (blow, move the tape, take a picture), keeps the pictures and a queue of
uploads in a local database, and serves a web page for controlling everything by hand.
The program also has a mock mode, so it can be developed on a laptop without the hardware.

\paragraph{Why a web page.}
Building and testing the device means trying things many times: moving the tape, switching
the fan, changing camera settings and taking pictures. A web page lets us do this from a laptop
or phone, without typing commands on the Pi. The main functions are:

\begin{itemize}
  \item \textbf{Hardware checks.} Checks that the camera is found and that wires are on the right
  pins, so building mistakes are found early.
  \item \textbf{Live preview and capture.} A live picture from the camera, and full-resolution
  pictures saved without compression, because compression destroys small details.
  \item \textbf{Camera settings.} Exposure, gain, white balance and image processing can be changed
  and are kept after a restart.
  \item \textbf{LED colour.} A colour wheel, an intensity slider and RGB values set the LED colour,
  which is also used for the pictures.
  \item \textbf{Manual control.} Fan on and off, moving the tape, and running a sampling cycle.
\end{itemize}

\paragraph{Notes and settings with every picture.}
We can write a note about each picture, and all camera and LED settings are saved with it.
Then we can go back and try settings that looked good when the setup or the place changes.

\fillin{Screenshot of the web page.}
\future{Upload of pictures to a central server, and a public map.}
